% =====================================================================
\begin{thebibliography}{99}
\addcontentsline{toc}{chapter}{Bibliografía}
\markboth{Bibliografía}{}
\small

\bibitem{bengio}
Y.~Bengio, R.~Ducharme, P.~Vincent y C.~Jauvin.
\newblock A neural probabilistic language model.
\newblock \emph{Journal of Machine Learning Research}, 3:1137--1155, 2003.
\newblock \url{https://www.jmlr.org/papers/v3/bengio03a.html}

\bibitem{bishopdl}
C.~M. Bishop y H.~Bishop.
\newblock \emph{Deep Learning: Foundations and Concepts}.
\newblock Springer, 2024.
\newblock Cap.~6, «Deep Neural Networks»; cap.~7, «Gradient Descent»; cap.~8, «Backpropagation».
\newblock \url{https://www.bishopbook.com/}

\bibitem{cmu11785}
B.~Raj \emph{et al.}
\newblock 11-785: Introduction to Deep Learning.
\newblock Carnegie Mellon University, otoño de 2025.
\newblock \url{https://deeplearning.cs.cmu.edu/}

\bibitem{cybenko}
G.~Cybenko.
\newblock Approximation by superpositions of a sigmoidal function.
\newblock \emph{Mathematics of Control, Signals and Systems}, 2(4):303--314, 1989.

\bibitem{glorot}
X.~Glorot e Y.~Bengio.
\newblock Understanding the difficulty of training deep feedforward neural networks.
\newblock En \emph{Proceedings of the Thirteenth International Conference on Artificial Intelligence and Statistics (AISTATS)}, PMLR~9:249--256, 2010.
\newblock \url{https://proceedings.mlr.press/v9/glorot10a.html}

\bibitem{gbc}
I.~Goodfellow, Y.~Bengio y A.~Courville.
\newblock \emph{Deep Learning}.
\newblock MIT Press, 2016.
\newblock Cap.~6, «Deep Feedforward Networks», y §8.4, «Parameter Initialization Strategies».
\newblock \url{https://www.deeplearningbook.org/}

\bibitem{he}
K.~He, X.~Zhang, S.~Ren y J.~Sun.
\newblock Delving deep into rectifiers: Surpassing human-level performance on ImageNet classification.
\newblock En \emph{Proceedings of the IEEE International Conference on Computer Vision (ICCV)}, 2015.
\newblock \url{https://arxiv.org/abs/1502.01852}

\bibitem{hendrycks}
D.~Hendrycks y K.~Gimpel.
\newblock Gaussian error linear units (GELUs).
\newblock Prepublicación arXiv:1606.08415, 2016.
\newblock \url{https://arxiv.org/abs/1606.08415}

\bibitem{hornik}
K.~Hornik.
\newblock Approximation capabilities of multilayer feedforward networks.
\newblock \emph{Neural Networks}, 4(2):251--257, 1991.

\bibitem{jm3}
D.~Jurafsky y J.~H. Martin.
\newblock \emph{Speech and Language Processing}, 3.ª ed., borrador del 19 de agosto de 2026.
\newblock Cap.~6, «Neural Networks».
\newblock \url{https://web.stanford.edu/~jurafsky/slp3/}

\bibitem{leshno}
M.~Leshno, V.~Y. Lin, A.~Pinkus y S.~Schocken.
\newblock Multilayer feedforward networks with a nonpolynomial activation function can approximate any function.
\newblock \emph{Neural Networks}, 6(6):861--867, 1993.

\bibitem{minsky}
M.~Minsky y S.~Papert.
\newblock \emph{Perceptrons: An Introduction to Computational Geometry}.
\newblock MIT Press, 1969.

\bibitem{novikoff}
A.~B.~J. Novikoff.
\newblock On convergence proofs on perceptrons.
\newblock En \emph{Proceedings of the Symposium on the Mathematical Theory of Automata}, vol.~12, págs.~615--622. Polytechnic Institute of Brooklyn, 1962.

\bibitem{radford2018}
A.~Radford, K.~Narasimhan, T.~Salimans e I.~Sutskever.
\newblock Improving language understanding by generative pre-training.
\newblock Informe técnico, OpenAI, 2018.
\newblock \url{https://cdn.openai.com/research-covers/language-unsupervised/language_understanding_paper.pdf}

\bibitem{radford2019}
A.~Radford, J.~Wu, R.~Child, D.~Luan, D.~Amodei e I.~Sutskever.
\newblock Language models are unsupervised multitask learners.
\newblock Informe técnico, OpenAI, 2019.
\newblock \url{https://cdn.openai.com/better-language-models/language_models_are_unsupervised_multitask_learners.pdf}

\bibitem{rosenblatt}
F.~Rosenblatt.
\newblock The perceptron: A probabilistic model for information storage and organization in the brain.
\newblock \emph{Psychological Review}, 65(6):386--408, 1958.

\bibitem{rumelhart}
D.~E. Rumelhart, G.~E. Hinton y R.~J. Williams.
\newblock Learning representations by back-propagating errors.
\newblock \emph{Nature}, 323:533--536, 1986.

\bibitem{shazeer}
N.~Shazeer.
\newblock GLU variants improve Transformer.
\newblock Prepublicación arXiv:2002.05202, 2020.
\newblock \url{https://arxiv.org/abs/2002.05202}

\bibitem{mit18065}
G.~Strang.
\newblock 18.065: Matrix Methods in Data Analysis, Signal Processing, and Machine Learning.
\newblock Massachusetts Institute of Technology, MIT OpenCourseWare, primavera de 2018.
\newblock Lección~27, «Backpropagation»; problemas de la sección VII.2 del libro \emph{Linear Algebra and Learning from Data}.
\newblock \url{https://ocw.mit.edu/courses/18-065-matrix-methods-in-data-analysis-signal-processing-and-machine-learning-spring-2018/}

\bibitem{telgarsky}
M.~Telgarsky.
\newblock Benefits of depth in neural networks.
\newblock En \emph{Conference on Learning Theory (COLT)}, 2016.
\newblock \url{https://arxiv.org/abs/1602.04485}

\bibitem{touvron}
H.~Touvron, T.~Lavril, G.~Izacard, X.~Martinet \emph{et al.}
\newblock LLaMA: Open and efficient foundation language models.
\newblock Prepublicación arXiv:2302.13971, 2023.
\newblock \url{https://arxiv.org/abs/2302.13971}

\bibitem{vaswani}
A.~Vaswani, N.~Shazeer, N.~Parmar, J.~Uszkoreit, L.~Jones, A.~N. Gomez, Ł.~Kaiser e I.~Polosukhin.
\newblock Attention is all you need.
\newblock En \emph{Advances in Neural Information Processing Systems 30}, 2017.
\newblock \url{https://arxiv.org/abs/1706.03762}

\bibitem{cs4780}
K.~Q. Weinberger.
\newblock CS4780/5780: Machine Learning for Intelligent Systems.
\newblock Cornell University, otoño de 2018. Lección~3, «The Perceptron».
\newblock \url{https://www.cs.cornell.edu/courses/cs4780/2018fa/lectures/lecturenote03.html}

\bibitem{cs224n}
D.~Yang e Y.~Choi.
\newblock CS224N: Natural Language Processing with Deep Learning.
\newblock Stanford University, invierno de 2026. Clase~3, «Backpropagation and Neural Network Basics», y clase~4, «Language Models and RNNs».
\newblock \url{https://web.stanford.edu/class/cs224n/}

\bibitem{cs224n2018}
C.~Manning y R.~Socher.
\newblock CS224N: Natural Language Processing with Deep Learning.
\newblock Stanford University, invierno de 2018. Tarea~1.
\newblock \url{https://web.stanford.edu/class/archive/cs/cs224n/cs224n.1184/assignment1/assignment1.pdf}

\bibitem{yang}
Z.~Yang, Z.~Dai, R.~Salakhutdinov y W.~W. Cohen.
\newblock Breaking the softmax bottleneck: A high-rank RNN language model.
\newblock En \emph{International Conference on Learning Representations (ICLR)}, 2018.
\newblock \url{https://arxiv.org/abs/1711.03953}

\bibitem{cs231n}
F.-F. Li \emph{et al.}
\newblock CS231n: Deep Learning for Computer Vision. Notas del curso.
\newblock Stanford University, 2026. «Backpropagation, Intuitions» y «Neural Networks», partes~1 y~2.
\newblock \url{https://cs231n.github.io/}

\bibitem{d2l}
A.~Zhang, Z.~C. Lipton, M.~Li y A.~J. Smola.
\newblock \emph{Dive into Deep Learning}, versión 1.0.3.
\newblock Cambridge University Press, 2023.
\newblock Cap.~5, «Multilayer Perceptrons», y §9.3, «Language Models».
\newblock \url{https://d2l.ai/}

\end{thebibliography}
